% STATUS: DRAFT
\chapter{Weyl relations and Stone--von Neumann}\label{ch:weyl}
\skippable

\begin{physmotiv}
The canonical commutation relations $[x,p]=\ii$ cannot hold for bounded operators (\cref{ch:hilbert}) and are only formally meaningful for unbounded ones. The cure is to exponentiate. The resulting \emph{Weyl relations} are a statement about unitaries only. Then comes the dramatic result: for finitely many degrees of freedom, there is essentially only one way to represent them (Stone--von Neumann), which is why textbook QM works; for infinitely many, there are uncountably many, which is why QFT needs algebras.
\end{physmotiv}

\section{The Weyl form of the CCR}

Define $U(a)=\ee^{\ii ap}$ and $V(b)=\ee^{\ii bx}$ for $a,b\in\R$. Using $\ee^{\ii ap}x\ee^{-\ii ap}=x+a$ one finds
\begin{equation}
U(a)V(b)=\ee^{\ii ab}\,V(b)U(a),\qquad U(a)U(a')=U(a+a'),\quad V(b)V(b')=V(b+b').\label{eq:weyl}
\end{equation}
Combining, $W(a,b)=\ee^{-\ii ab/2}U(a)V(b)$ satisfies $W(z)W(z')=\ee^{\ii\sigma(z,z')/2}W(z+z')$ with symplectic form $\sigma(z,z')=ab'-ba'$ for $z=(a,b)$, $z'=(a',b')$. More generally for $n$ degrees of freedom, or a field, $W(f)$ labelled by a vector $f$ in a real symplectic space $(\mathcal S,\sigma)$:
\begin{equation}
W(f)W(g)=\ee^{\frac\ii2\sigma(f,g)}W(f+g),\qquad W(f)\ad=W(-f).
\end{equation}
For the free scalar field, $\mathcal S$ is a space of test-function data, $W(f)=\ee^{\ii\phi(f)}$ with $\phi(f)$ the smeared field, and $\sigma(f,g)$ is the commutator function (the symplectic form of the classical phase space). These $W(f)$ generate the \emph{Weyl algebra} (a \cstar-algebra, \cref{ch:cstar}), which is the prototypical algebra for bosonic QFT. Everything is bounded: $W(f)$ is unitary.

\section{Stone--von Neumann}

A representation of \eqref{eq:weyl} on $\HH$ is \emph{regular} if $U(a)$, $V(b)$ are strongly continuous in $a,b$ (so that $x,p$ exist as self-adjoint generators by Stone's theorem).

\begin{theorem}[Stone--von Neumann]
For $n<\infty$ degrees of freedom, every regular irreducible representation of the Weyl relations is unitarily equivalent to the Schr\"odinger representation on $L^2(\R^n)$. Every regular representation is a direct sum of copies of it.
\end{theorem}

This is why finite-dimensional QM has no representation ambiguity: ``the'' Hilbert space is $L^2(\R^n)$, Fock-space and Schr\"odinger descriptions of the harmonic oscillator are unitarily equivalent, and so are different choices of $\omega$ in the Fock vacuum.

\begin{whatwrong}
Regularity is essential. The Weyl relations also have non-regular representations (e.g.\ on $\ell^2(\R_{\rm discrete})$ where $U(a)$ translates a discrete basis and $V(b)$ is multiplication, with no strongly continuous $V$) in which there is no $x$ operator. Such representations appear in polymer quantization (loop quantum gravity); they are inequivalent to Schr\"odinger. Physics selects regular ones because we want $x,p$ to exist.
\end{whatwrong}

\section{The failure for infinitely many degrees of freedom}

For a free field on a Cauchy surface, the Fock vacuum $\Omega_m$ of mass $m$ defines a representation of the Weyl algebra on the Fock space $\mathcal F_m$. The two vacua $\Omega_m,\Omega_{m'}$ are related by a Bogoliubov transformation with coefficients $\beta_k\simeq\frac{m'^2-m^2}{4\omega_k^{2}}$ at large momentum, one for each mode $k$. A Bogoliubov transformation is implementable by a unitary operator iff $\sum_k|\beta_k|^2<\infty$ (Shale's criterion). In a finite box the sum over modes converges (in $3+1$ dimensions $\int\dd^3k\,k^{-4}<\infty$), so the two representations are equivalent. In \emph{infinite volume}, the sum becomes $V\int\frac{\dd^3k}{(2\pi)^3}|\beta_k|^2$ with $V\to\infty$, and diverges: the two representations are unitarily \emph{inequivalent}. This is the same phenomenon as the spin chain of \cref{ch:why_algebras}: in infinite volume two different states differ by infinitely many excitations, so no vector of one Hilbert space represents a state of the other (equivalently, the overlap of the two vacua vanishes: see the last exercise).

\emph{(The ultraviolet is a separate issue: the interacting vacuum is not equivalent to a free one even in finite volume, which is Haag's theorem.)}

\begin{keyidea}
For finitely many degrees of freedom, the algebra determines the representation (S--vN). For infinitely many, it does not; different physical situations (masses, temperatures, densities) live in inequivalent representations of one and the same Weyl algebra. This is the reason the algebra, not $\HH$, is fundamental.
\end{keyidea}

We return to this in \cref{ch:inequiv}, where the thermal representation of the Weyl algebra at temperature $T>0$ is shown to be inequivalent to the vacuum representation.

\section*{Exercises}
\begin{exercise}
Verify $\ee^{\ii ap}x\ee^{-\ii ap}=x+a$ (use $\frac{\dd}{\dd a}$) and derive \eqref{eq:weyl}.
\end{exercise}
\begin{exercise}
Show that if $[x,p]=\ii$ holds on a finite-dimensional space it would imply $\Tr\ii\id=0$, so $\dim\HH=\infty$. Compare with the Wintner--Wielandt argument.
\end{exercise}
\begin{exercise}
Quantum mechanics on a circle: on $L^2(S^1)$ the operator $x$ is not well defined, but $V(n)=\ee^{\ii nx}$ is for integer $n$. Show that $U(a)$ (translations) and $V(n)$ satisfy the Weyl relation only for integer $n$. Show that replacing $p\to p+\theta$ gives a family of inequivalent representations of this smaller algebra, labelled by $\theta\in[0,1)$. How does this differ from the line, where Stone--von Neumann forbids such a family?
\end{exercise}
\begin{exercise}
Two oscillator vacua with frequencies $\omega,\omega'$ are related by a Bogoliubov transformation with $\alpha=\frac{\omega+\omega'}{2\sqrt{\omega\omega'}}$, $\beta=\frac{\omega-\omega'}{2\sqrt{\omega\omega'}}$ (check $\alpha^2-\beta^2=1$). Compute (e.g.\ with Gaussians) the overlap $|\braket{0_\omega}{0_{\omega'}}|=\alpha^{-1/2}=(1+\beta^2)^{-1/4}$. For $N$ independent oscillators show the total overlap is $\prod_j(1+\beta_j^2)^{-1/4}$ and that it has a non-zero limit as $N\to\infty$ iff $\sum_j\beta_j^2<\infty$.
\end{exercise}
